\subsection{西群与正交群。}

设 \(\Phi\) 是 E 上的一个埃尔米特型；A-模 E 中保持 \(\Phi\) 不变的自同构称为关于 \(\Phi\) 的西自同构（或西变换），它们所成的群称为与 \(\Phi\) 相伴的西群；记作 \(\mathbf{U}(\Phi)\) 。给定 E 上的一个二次型 \(Q \neq 0\) ，A-模 E 中保持 Q 不变的自同构称为关于 Q 的正交自同构（或正交变换）；它们所成的群称为与 Q 相伴的正交群；记作 \(\mathbf{O}(Q)\) 。

对二次型 Q 的每个正交变换都是对与 Q 相伴的双线性型的西变换。当数量 2 在 A 中不等于 0 也不是零因子时（§ 3，第 4 小节，(13)），逆命题成立，例如当 A 是特征 ≠ 2 的体时。

特别地，考虑模 \(E = A^{n}\) 上的埃尔米特型 \(\Phi_{0}\) ，它关于 E 的典范基 \((e_{i})\) 的矩阵是单位矩阵 \(I_{n}\) 。与 \(\Phi_{0}\) 相伴的西自同构简称为 n 变量的西自同构（或西变换）；它们所成的群称为 n 变量西群，有时记作 \(\mathbf{U}(n, \mathbf{A})\) 或 \(\mathbf{U}_{n}(\mathbf{A})\) 。西自同构关于 \((e_{i})\) 的矩阵 U 称为酉矩阵。这样的矩阵是可逆的，并且依 § 1 第 10 小节的公式 (48)，满足关系

\[
{ } ^ { t } U . \overline { { U } } = I _ { n } ;\tag{4}
\]

反之，若 A 是一个交换环或一个体，则满足 (4) 的矩阵 U 是可逆的，从而是酉的。

当 J 是恒等映射且 2 在 A 中既不等于 0 也不是零因子时，我们用 n 变量正交群、n 变量正交自同构（或正交变换）以及正交矩阵这些术语来代替前面的术语，并把 \(\mathbf{O}(n,\mathrm{A})\) （相应地，\(\mathbf{O}_n(\mathrm{A})\) ）用来代替 \(\mathbf{U}(n,\mathrm{A})\) （相应地，\(\mathbf{U}_n(\mathrm{A})\) ）。此时关系 (4) 写成

\[
{ } ^ { t } U . U = I _ { n }\tag{5}
\]

并且由于 A 是交换的，它是 U 为正交矩阵的一个必要充分条件。

命题 3. --- 假设 A 是一个交换域且 E 的维数有限 \textgreater{} 0。设 Φ 是 E 上的一个非退化埃尔米特型。映射 u → det u 是从与 Φ 相伴的酉群 U(Φ) 到 A 的乘法子群 H 上的一个同态，H 由使 ρρ̄ = 1 的元素 ρ 组成（当 J 是恒等映射时，这是化为 \{1, -1\} 的子群）。

实际上，设 \(u\) 是 \(\mathbf{U}(\Phi)\) 的一个元素，\(U\) 是它关于 E 的一个基的矩阵，\(R\) 是 \(\Phi\) 关于这个基的矩阵。关系 \(R = {}^t U.R.\overline{U}\) ( \(\S 1\) ，第 10 小节，公式 (48)) 表明有 \((\det U)(\det \overline{U}) = 1\) ，因为 \(R\) 可逆；由此得 \((\det u)(\overline{\det u}) = 1\) 。同态 \(u \to \det u\) 把 \(\mathbf{U}(\Phi)\) 映到 H 上。实际上，当 A 的特征是 2 且 J 是恒等映射时，H 化为元素 1。否则存在 E 的一个正交基 \((e_i)(i = 1,\ldots,n)\) （定理 1）；对每个 \(\rho \in A\) （使 \(\rho \bar{\rho} = 1\) ），设 \(u\) 是 E 的由 \(u(e_1) = \rho e_1\) 与 \(u(e_i) = e_i\) （对 \(i = 2,\ldots,n\) ）定义的自同构；则 \(u\) 是酉的且 det \(u = \rho\) ，由此即得命题。

在命题 3 的条件下，同态 \(u \to \det u\) 的核是 \(\mathbf{U}(\Phi)\) 的一个正规子群，称为与 \(\Phi\) 相伴的特殊酉群；有时记作 \(\mathbf{SU}(\Phi)\) 。

当 J 是恒等映射且 A 的特征不是 2 时，这个群也称为与 \(\Phi\) （或与二次型 \(Q(x) = \Phi(x, x)\) ）相伴的特殊正交群，有时记作 SO(Q)。

若 E = A \(^{n}\) 且 Φ 是关于 E 的典范基的矩阵为单位矩阵的那个型，则使用记号 SU(n, A) 或 SU \(_{n}\) (A) 与 SO(n, A) 或 SO \(_{n}\) (A)。

